% SHAWAF Stage 1 — technical report (compile with pdflatex or latexmk)
% Figures: optional \includegraphics from ../assets/ when building from docs/
\documentclass[11pt,a4paper]{article}
\usepackage[margin=2.5cm]{geometry}
\usepackage{amsmath,amssymb}
\usepackage{booktabs}
\usepackage{graphicx}
\usepackage{hyperref}
\usepackage{enumitem}
\usepackage{listings}
\usepackage{xcolor}
\usepackage{longtable}
\usepackage{caption}

\hypersetup{colorlinks=true, linkcolor=blue, urlcolor=blue, citecolor=blue}
\lstset{
  basicstyle=\ttfamily\small,
  breaklines=true,
  frame=single,
  backgroundcolor=\color{gray!8},
}

\title{SHAWAF Stage 1:\\Text-to-Tracklet Retrieval with Frozen Vision--Language Encoders\\[0.4em]
\large Sampling, Pooling (Method~C), Benchmarks, and Interactive Retrieval}
\author{Benchmarking Video/Image Language Models for Tracklet Retrieval}
\date{\today}

\begin{document}
\maketitle

\begin{abstract}
SHAWAF (Stage~1) evaluates \emph{frozen} vision--language models (VLMs) for natural-language query $\rightarrow$ tracklet retrieval: each text description is embedded, each gallery tracklet is embedded from video frames or still images, and cosine similarity ranks the gallery.
This report documents design choices in the \texttt{Embeddings} repository: Tier-A datasets, encoder groups (\texttt{base} vs.\ \texttt{extended\_modern} vs.\ \texttt{extended\_gme}), frame sampling and sliding-window protocols, pooling rules including \textbf{Type~C} (\texttt{global\_local} hybrid scoring), benchmark results, the numpy/Qdrant indexing path, the Gradio demo, limitations, and a roadmap toward large-scale corpus ingestion and optional two-stage temporal grounding.
All headline benchmark numbers below are taken from the project README unless your in-flight Kaggle runs supersede them; pooling comparisons with \texttt{global\_local} should be read from fresh notebook outputs.
\end{abstract}

\tableofcontents
\newpage

%=============================================================================
\section{Introduction and problem statement}
%=============================================================================

\subsection{What we are solving}
Given a \textbf{gallery} of tracklets (short person-centric video clips or still-image identities) and a set of \textbf{text queries} (captions describing appearance, and sometimes motion or scene), we rank gallery items by semantic similarity to each query.
Positives share the same \textbf{person identity} as the query's source item; metrics are standard re-identification retrieval: Rank-1, mean average precision (mAP), and related variants.

This is \textbf{text-to-tracklet retrieval}, not:
\begin{itemize}[noitemsep]
  \item classical image-to-image Re-ID (no text query);
  \item open-world video question answering;
  \item pixel-level referring video segmentation (RVOS);
  \item precise \textbf{temporal grounding} inside one long untrimmed video (start/end timestamps), though that is discussed as future work.
\end{itemize}

\subsection{Why frozen encoders only (Stage~1)}
Stage~1 deliberately avoids fine-tuning on target benchmarks (except where a \emph{pretrained} checkpoint was already task-specific, e.g.\ IRRA).
The scientific question is: \emph{which off-the-shelf encoder and which temporal pooling protocol} best align text with tracklets under zero-shot evaluation.
Later SHAWAF stages may fuse Re-ID features or train projections; they are out of scope here.

\subsection{Relation to ``hours of video in a database''}
The \textbf{product-shaped} goal---ingest many hours of footage, index embeddings, return top-$K$ clips for a text query---uses the \textbf{same retrieval stack} as the benchmark (encode chunks $\rightarrow$ vector store $\rightarrow$ nearest-neighbor search).
Tier-A benchmarks use \textbf{pre-cut tracklets}; scaling to raw long video requires an explicit \textbf{chunking} policy (sliding windows or detector/tracklet pipelines) documented in Section~\ref{sec:corpus}.
That is engineering scale-up, not a different problem family from Stage~1 retrieval.

%=============================================================================
\section{Repository architecture}
%=============================================================================

The Python package \texttt{shawaf\_vlm/} implements loaders, encoders, sampling, pooling, metrics, and evaluation (\texttt{eval\_loop.py}).
Scripts under \texttt{scripts/} download Tier-A data, run indexing, and smoke-test the demo.
\texttt{demo/app.py} is a thin Gradio UI over \texttt{shawaf\_vlm.retrieval}.
Notebooks (\texttt{kaggle-zero-shot.ipynb}, \texttt{colab\_extended\_zero\_shot.ipynb}, \texttt{pooling\_methods\_comparison.ipynb}) reproduce large benchmark grids and pooling ablations.

\begin{description}
  \item[Benchmark path] Dataset splits $\rightarrow$ frame/window extraction $\rightarrow$ \texttt{encode\_videos} / \texttt{encode\_texts} $\rightarrow$ pool $\rightarrow$ similarity matrix $\rightarrow$ Rank-1 / mAP.
  \item[Interactive path] \texttt{download\_tier\_a.py} $\rightarrow$ \texttt{index\_qdrant.py} $\rightarrow$ on-disk numpy manifest (+ optional Qdrant) $\rightarrow$ demo search with \texttt{mean}, \texttt{query\_max}, or \texttt{global\_local}.
\end{description}

%=============================================================================
\section{Datasets (Tier A and beyond)}
%=============================================================================

\label{sec:datasets}

Table~\ref{tab:datasets} summarizes datasets integrated for Stage~1.
All listed Tier-A sets are used \textbf{zero-shot} for every general VLM; IRRA is task-trained on person retrieval (not on these exact splits), which partially explains its strength.

\begin{table}[ht]
\centering
\caption{Tier-A text-to-tracklet benchmarks in SHAWAF.}
\label{tab:datasets}
\begin{tabular}{@{}llll@{}}
\toprule
Dataset & Modality & Gallery scale (test) & Query style \\
\midrule
TVPReid PRID & Video tracklets & 185 & 2 captions / video \\
TVPReid iLIDS & Video tracklets & 75 & 2 captions / video \\
TVPReid Duke & Video tracklets & 603 & 2 captions / video \\
RSTPReid & Still images (5 cams) & 1{,}000 images & 2 captions / image \\
GroOT-MOT17 & MOT17 track mp4s & 454 tracks & appearance / action / combined \\
\bottomrule
\end{tabular}
\end{table}

\textbf{TVPReid} captions emphasize \textbf{appearance} with some motion/scene wording; evaluation is video-level (retrieve the source mp4).
\textbf{RSTPReid} removes temporal signal (single frame repeated to fill clip encoders).
\textbf{GroOT-MOT17} separates \textbf{appearance}, \textbf{action}, and \textbf{combined} queries on the same tracks; action captions are highly repetitive (dataset ceiling; see Section~\ref{sec:limitations}).

\textbf{ActivityNet Captions} (1{,}000-video sample) is an out-of-domain long-video check in the README; ranking \emph{reverses} relative to person retrieval (video-pretrained models lead).

\textbf{Planned / stub:} TV-MARS, CUHK-PEDES, ICFG-PEDES (\texttt{docs/datasets\_catalog.md}).
\textbf{Out of scope:} RefAVA/RAVAR (classification), raw AVA without retrieval labels.

%=============================================================================
\section{Model selection and environment groups}
%=============================================================================

\subsection{Encoder registry}
Models are registered in \texttt{shawaf\_vlm/models/registry.py} with a common \texttt{VideoTextEncoder} protocol: \texttt{encode\_videos(list[frames])} and \texttt{encode\_texts(list[str])}, L2-normalized embeddings, cosine similarity at evaluation time.

\subsection{Why three Kaggle/Colab groups}
Dependency profiles are \textbf{not mutually compatible} at a single \texttt{transformers} version:
\begin{itemize}[noitemsep]
  \item \textbf{base} --- SigLIP~2, Perception Encoder (PE-Core L/14), IRRA, InternVideo2 CLIP-S, InternVideo2-1B-s2, X-CLIP, LanguageBind.
  \item \textbf{extended\_modern} --- OpenAI CLIP ViT-L/14, Jina CLIP v2, Qwen3-VL-Embedding-2B (\texttt{transformers==4.57.3}).
  \item \textbf{extended\_gme} --- GME-Qwen2-VL-2B (\texttt{transformers==4.51.3}).
\end{itemize}
Each group is typically one GPU session on Kaggle; results merge via saved \texttt{zero\_shot\_results/} folders.

\subsection{Rationale for including each family}
\begin{description}
  \item[IRRA] Task-specific text--person alignment; upper bound for appearance descriptions on TVPReid/RSTP.
  \item[SigLIP / PE-Core] Strong general image--language baselines; competitive on appearance and ActivityNet sentences.
  \item[InternVideo2 CLIP-S] Video-language pretraining; strong on long events and GroOT combined captions.
  \item[Qwen3-VL-Embedding] Best among \emph{extended} general models on TVPReid; heavier inference.
  \item[X-CLIP / LanguageBind] Legacy video--text baselines; useful cost--accuracy anchors.
  \item[InternVideo2-1B-s2] Short 4-frame native clip; included for completeness, weak on retrieval here.
\end{description}

Native clip lengths differ (e.g.\ 8 frames vs.\ 4 for s2); protocols respect \texttt{NATIVE\_FRAMES} in notebooks.

%=============================================================================
\section{Frame sampling and video decoding}
%=============================================================================

\subsection{Uniform clip sampling (\texttt{uniform8})}
For a single clip per tracklet, \texttt{sample\_frame\_paths} or \texttt{resolve\_frame\_paths} selects \texttt{num\_frames} (default 8) indices along the tracklet.
If the tracklet is shorter, frames are \textbf{padded by repeating the last frame} so encoders always see fixed length.
For a single \texttt{.mp4}, frames are decoded with OpenCV or torchvision, written as cached JPEGs under \texttt{.cache/frame\_cache} (or per-video \texttt{.frames\_N}).

Sampling modes:
\begin{itemize}[noitemsep]
  \item \textbf{uniform} --- linspace over frame indices;
  \item \textbf{middle} --- midpoint of each temporal bin (InternVideo2-style).
\end{itemize}

\subsection{Dense sampling for sliding windows}
Protocol \texttt{reid\_8fps\_n64} (and variants \texttt{vt\_1fps\_n12}, \texttt{vt\_2fps\_n32}) first builds a \textbf{dense} frame list:
\texttt{resolve\_dense\_frame\_paths} samples video at \texttt{sample\_fps} (e.g.\ 8\,fps) up to \texttt{max\_frames} (e.g.\ 64), covering the whole tracklet.
Still images use subsampled crop lists capped at \texttt{max\_frames}.

\subsection{Why cache frames to JPEG}
Most encoders consume \textbf{image paths}, not raw tensors in the eval loop.
Caching avoids re-decoding mp4v/MPEG-4 clips on every run (TVPReid tracklets often use \texttt{mp4v}, which browsers also do not play without transcoding).

%=============================================================================
\section{Sliding windows and clip splitting}
%=============================================================================

\label{sec:windows}

Given dense frames $f_1,\ldots,f_T$, \texttt{sliding\_windows} produces overlapping clips of length $W$ (default $W{=}8$) with stride $S$ (default $S{=}4$):
\[
\mathcal{C} = \bigl\{ (f_i,\ldots,f_{i+W-1}) \;:\; i \in \{0,S,2S,\ldots\} \cup \{\text{last valid start}\} \bigr\}.
\]
If $T \le W$, a single clip is formed via uniform resampling to $W$ frames.

\textbf{Encode once, pool many:} All windows are encoded in one batched \texttt{encode\_videos} pass; features are grouped per tracklet for pooling (Section~\ref{sec:pooling}).
This is the main \textbf{weak temporal} mechanism in Stage~1: temporal extent enters only through which frames appear in which window, not through a recurrent or transformer-over-time aggregator inside the benchmark loop.

\subsection{Protocol summary}
\begin{table}[ht]
\centering
\caption{Common temporal protocols (see README for headline selection per model).}
\begin{tabular}{@{}llll@{}}
\toprule
Name & Dense sampling & Window $W$ & Stride $S$ \\
\midrule
\texttt{uniform8} & N/A (one clip) & 8 & --- \\
\texttt{vt\_1fps\_n12} & 1\,fps, max 12 & 8 & 4 \\
\texttt{vt\_2fps\_n32} & 2\,fps, max 32 & 8 & 4 \\
\texttt{reid\_8fps\_n64} & 8\,fps, max 64 & 8 & 4 \\
\bottomrule
\end{tabular}
\end{table}

%=============================================================================
\section{Pooling and scoring (including Method C)}
%=============================================================================

\label{sec:pooling}

For each tracklet $g$, let clip embeddings be $\mathbf{c}_1,\ldots,\mathbf{c}_K \in \mathbb{R}^D$ (L2-normalized).
Let query embedding be $\mathbf{q}$ (L2-normalized).

\begin{description}
  \item[\texttt{mean}] Global vector $\bar{\mathbf{c}} = \mathrm{normalize}(\frac{1}{K}\sum_k \mathbf{c}_k)$; score $s = \mathbf{q}^\top \bar{\mathbf{c}}$.
  \item[\texttt{max}] Element-wise max over clips then normalize (CLIP4Clip-style); less used in headline tables.
  \item[\texttt{mean\_s8}] Subsample windows every 8 frames (approx.\ stride-8 from stride-4 encoding) before mean pool.
  \item[\texttt{query\_max}] \textbf{Local / late interaction:} $s = \max_k \mathbf{q}^\top \mathbf{c}_k$.
  \item[\texttt{global\_local} (Type C / Method C)] Hybrid:
  \[
  s = \alpha\, (\mathbf{q}^\top \bar{\mathbf{c}}) + (1-\alpha)\, \max_k (\mathbf{q}^\top \mathbf{c}_k),
  \]
  with default $\alpha = 0.5$ (\texttt{HYBRID\_ALPHA} in notebooks/CLI).
\end{description}

\textbf{Why Type C?}
Mean pooling is stable when the query describes the whole tracklet (paragraph-level, overall appearance).
\texttt{query\_max} helps when only \textbf{one temporal segment} matches the text (event-like sentence).
Type C interpolates without re-encoding clips---all pools share the same encoded features in \texttt{evaluate\_text\_to\_tracklet\_windows}.

\textbf{Comparison runs:}
Use \texttt{pooling\_methods\_comparison.ipynb} for one model/dataset, or \texttt{kaggle-zero-shot.ipynb} section~5.3b after a full eval with \texttt{POOLS} containing \texttt{global\_local}.

%=============================================================================
\section{Evaluation metrics and protocol}
%=============================================================================

\begin{itemize}[noitemsep]
  \item Embeddings L2-normalized; distance $= 1 - \cosine$.
  \item Rank-1 and mAP with identity labels; optional same-camera junk filtering (\texttt{junk\_same\_camera}).
  \item Headline README numbers pick the \textbf{best} tested protocol and pool per model/dataset (not a single fixed setting for all rows).
  \item GroOT \textbf{action} column often excluded from cost--accuracy plots due to caption repetition ceiling.
\end{itemize}

%=============================================================================
\section{Benchmark results (reference)}
%=============================================================================

\label{sec:results}

The following tables are copied from the project README for writing convenience.
\textbf{Update} with your new Kaggle \texttt{base} + \texttt{extended\_modern} runs when they finish, especially rows that add \texttt{global\_local} comparisons.

\subsection{TVPReid (Rank-1 / mAP \%)}
\begin{table}[ht]
\centering
\small
\begin{tabular}{@{}lrrr@{}}
\toprule
Model & PRID & iLIDS & Duke \\
\midrule
IRRA & 66.55 / 76.09 & 30.00 / 41.29 & 38.06 / 50.79 \\
SigLIP 2 & 35.56 / 50.84 & 18.67 / 28.88 & 21.97 / 33.03 \\
PE-Core L/14 & 38.38 / 52.04 & 19.33 / 32.55 & 17.66 / 28.12 \\
InternVideo2 CLIP-S & 32.39 / 43.77 & 13.33 / 24.14 & 18.49 / 28.36 \\
Qwen3-VL-Emb-2B & 24.65 / 37.00 & 11.33 / 21.37 & 13.18 / 22.73 \\
\bottomrule
\end{tabular}
\end{table}

\subsection{RSTPReid (Rank-1 / mAP \%)}
IRRA 53.30 / 39.71; SigLIP 2 41.00 / 29.34; PE-Core 37.40 / 29.09 (see README for full table).

\subsection{GroOT-MOT17 (Rank-1 \%)}
Appearance / action / combined columns show appearance models win on appearance; combined benefits PE-Core and IRRA; action remains low (2--3\%) due to caption duplication.

\subsection{ActivityNet sample (out of domain)}
InternVideo2 CLIP-S leads paragraph; PE-Core leads sentence mean Rank-1; IRRA is weak---person specialization does not transfer to general events.

%=============================================================================
\section{Interactive retrieval stack}
%=============================================================================

\subsection{Indexing}
\texttt{scripts/index\_qdrant.py} encodes the \textbf{gallery only} with the same windowing as eval (\texttt{encode\_gallery\_tracklets}), then writes:
\begin{itemize}[noitemsep]
  \item \texttt{manifest.json} + \texttt{vectors.npz} under \texttt{\$HOME/.cache/shawaf\_vlm/indexes/<model>/<dataset>/};
  \item optional Qdrant collection (\texttt{docker compose up qdrant}) unless \texttt{--skip-qdrant}.
\end{itemize}
Each tracklet stores \texttt{global\_vector}, all \texttt{clip\_vectors}, and \texttt{media\_path} for preview.

\subsection{Search (demo)}
\texttt{demo/app.py} encodes the user query once, scores all indexed tracklets with selected mode (\texttt{mean}, \texttt{query\_max}, \texttt{global\_local}), returns top-5.
Previews transcode \texttt{mp4v} to H.264 via \texttt{imageio-ffmpeg} when available.

\subsection{What the demo is not}
It does not run full benchmark mAP; it does not guarantee action understanding for free-form queries (``woman drinking water'') beyond what the frozen encoder provides.

%=============================================================================
\section{Corpus-scale ingestion (roadmap)}
%=============================================================================

\label{sec:corpus}

To index \textbf{hours} of raw video:
\begin{enumerate}
  \item \textbf{Chunk} --- sliding windows on each file (reuse \texttt{sample\_fps}, \texttt{max\_frames}, \texttt{stride}) \emph{or} person tracklets from a detector/MOT pipeline.
  \item \textbf{Metadata} --- store \texttt{video\_id}, \texttt{t\_start}, \texttt{t\_end}, path per row (extend manifest schema).
  \item \textbf{Batch index} --- planned \texttt{index\_corpus.py} over a directory tree; same Qdrant/numpy search.
  \item \textbf{Query} --- identical vector top-$K$; UI shows time range + file.
\end{enumerate}

Ingest cost is paid \textbf{offline}; online search stays fast (milliseconds--hundreds of ms at scale with Qdrant).

%=============================================================================
\section{Coarse retrieval and temporal grounding (future)}
%=============================================================================

\label{sec:grounding}

\textbf{Coarse stage (Stage~1 + corpus index):} Text embedding vs.\ precomputed window/tracklet embeddings $\rightarrow$ top-$K$ candidates across the archive.
Fast, amortized indexing.

\textbf{Fine stage (optional research):} Temporal video grounding (TVG) or referring VOS on \textbf{one candidate clip} to predict $[t_{\mathrm{start}}, t_{\mathrm{end}}]$ or masks.
LVLM grounding (TimeChat, VTimeLLM, UniTime, Time-R1, ED-VTG, etc.) is \textbf{slow per clip} (often seconds on GPU) and should run on top-5 or on user request, not on every query.

This is \textbf{asymmetric retrieval}: wide cheap search, narrow expensive refinement---standard for large video archives.

SHAWAF Stage~1 does \textbf{not} implement Stage~B; it documents the interface (returned clip path + optional future timestamp fields).

%=============================================================================
\section{Limitations}
%=============================================================================

\label{sec:limitations}

\begin{enumerate}
  \item \textbf{Zero-shot general VLMs} are not trained on your ad-hoc queries; TVPReid-style appearance text works best.
  \item \textbf{Weak temporal modeling} --- windowing + max/hybrid pool, not explicit long-range video reasoning.
  \item \textbf{Action queries} (GroOT action, ``drinking water'') hit encoder and caption ceilings; repetitive action labels limit mAP interpretation.
  \item \textbf{Headline table heterogeneity} --- best protocol/pool per model; not apples-to-apples without fixed \texttt{reid\_8fps\_n64} + single pool.
  \item \textbf{ActivityNet numbers} are a 1{,}000-video sample, not official leaderboard scores.
  \item \textbf{mp4v} source clips require transcoding for in-browser preview; retrieval is unaffected.
  \item \textbf{IRRA} is not zero-shot in the strictest sense (person-retrieval pretraining); included as a strong specialist baseline.
  \item \textbf{No fusion} with dedicated Re-ID CNN embeddings in Stage~1.
\end{enumerate}

%=============================================================================
\section{Reproducibility}
%=============================================================================

\begin{itemize}[noitemsep]
  \item Install: \texttt{pip install -e ".[all]"} or model extras; demo: \texttt{pip install -e ".[demo]"} + \texttt{imageio-ffmpeg}.
  \item Data: \texttt{python scripts/download\_tier\_a.py}; set \texttt{HF\_TOKEN} if needed.
  \item Eval CLI: \texttt{shawaf-vlm-eval} / \texttt{scripts/eval\_encoder.py} with \texttt{--pools mean,query\_max,global\_local --hybrid-alpha 0.5}.
  \item Kaggle: \texttt{notebooks/kaggle-zero-shot.ipynb}, secret \texttt{hugging\_face} for gated checkpoints.
  \item Tests: \texttt{pytest} from repo root (\texttt{tests/test\_tv\_mars.py}, \texttt{test\_retrieval\_cache.py}, \texttt{test\_demo\_app.py}, etc.).
  \item Environments: \texttt{docs/ENVIRONMENTS.md} (Gradio/FastAPI pins, Hub version conflicts).
\end{itemize}

%=============================================================================
\section{Conclusion}
%=============================================================================

SHAWAF Stage~1 establishes a reproducible benchmark for frozen text-to-tracklet retrieval across TVPReid, RSTPReid, and GroOT-MOT17, with explicit control over \textbf{sampling}, \textbf{sliding windows}, and \textbf{pooling}, including Type~C \texttt{global\_local} fusion.
The same embedding pipeline powers an interactive demo and future corpus indexing.
Specialized models (IRRA) dominate person-description tasks; video-language models excel on general events; Type~C pooling tests whether combining global and local scores improves over either alone---your current Kaggle runs supply the definitive pooling comparison tables.
Next engineering steps are corpus chunking metadata and optional on-demand temporal grounding on retrieved clips, without changing the core Stage~1 metric definitions.

\appendix
\section{Quick reference: CLI pooling flags}
\begin{lstlisting}
python scripts/eval_encoder.py --model pe_core_l14 \
  --windows --stride 4 --sample-fps 8 --max-frames 64 \
  --pools mean,query_max,global_local --hybrid-alpha 0.5
\end{lstlisting}

\section{File index}
\begin{tabular}{@{}ll@{}}
\toprule
Path & Role \\
\midrule
\texttt{shawaf\_vlm/sampling.py} & Frames, FPS sampling, sliding windows \\
\texttt{shawaf\_vlm/pooling.py} & Mean, max, query\_max, global\_local \\
\texttt{shawaf\_vlm/eval\_loop.py} & Windowed encode + multi-pool metrics \\
\texttt{shawaf\_vlm/retrieval/} & Index, search, Qdrant \\
\texttt{demo/app.py} & Gradio UI \\
\texttt{docs/ARCHITECTURE.md} & Short architecture summary \\
\bottomrule
\end{tabular}

\end{document}
